\chapter{Couche limite turbulente : la méthode d'entraînement de Head}\label{ch:head}
\iterbanner{V0}

\section{Pourquoi une deuxième équation}
En laminaire, la forme du profil de vitesse suit le gradient de pression local, et un seul paramètre ($\lambda$) suffit. En turbulent, ce n'est plus vrai : la forme moyenne dépend de l'histoire de la couche, de la façon dont la turbulence s'est développée en amont. Il faut donc suivre $H$ comme une inconnue à part entière, et lui donner sa propre équation d'évolution.

\section{L'idée de Head : l'entraînement}
Le débit volumique qui traverse la couche limite, par unité d'envergure, vaut
\[
Q=\int_0^\delta u\,\dd y=U_e\,(\delta-\delta^*)=U_e\,\theta\,H_1,\qquad H_1=\frac{\delta-\delta^*}{\theta}.
\]
Ce débit ne peut augmenter que si la couche \textbf{entraîne} du fluide extérieur irrotationnel à travers sa frontière, l'interface turbulent--non turbulent (figure~\ref{fig:entrainement}). Head (1958) suppose que la vitesse d'entraînement, rapportée à $U_e$, ne dépend que de la forme du profil :
\[
\frac{1}{U_e}\frac{\dd Q}{\dd s}=\frac{1}{U_e}\frac{\dd(U_e\theta H_1)}{\dd s}=F(H_1),
\]
et que $H_1$ est lui-même une fonction de $H$. Avec l'équation de von Kármán, on obtient un système fermé de deux EDO.

\begin{figure}[htbp]
\centering
\begin{tikzpicture}[x=1.1cm,y=1.1cm]
\fill[pattern=north east lines,pattern color=cmuted] (0,-0.25) rectangle (10,0);
\draw[cink,thick] (0,0) -- (10,0);
\draw[cblue,very thick,decorate,decoration={snake,amplitude=0.6mm,segment length=4mm}] (0,0.8) .. controls (4,1.3) and (7,1.7) .. (10,2.1);
\node[cblue,font=\small] at (8.2,2.25) {frontière turbulente $\delta(s)$};
\foreach \x in {2.5,5.5,8}{\draw[corange,thick,-{Stealth[length=2mm]}] (\x-0.3,{0.8+0.13*\x+0.7}) -- (\x,{0.8+0.13*\x+0.12});}
\node[corange,font=\small] at (3.2,2.4) {entraînement de fluide extérieur};
\draw[cink,-{Stealth[length=2.5mm]},thick] (-1.2,1.0) -- (-0.2,1.0) node[midway,above,font=\small] {$U_e$};
\node[font=\small,align=center] at (5,0.55) {débit $Q=U_e(\delta-\delta^*)$ croissant};
\end{tikzpicture}
\caption{L'idée de Head : le débit dans la couche limite turbulente n'augmente que par entraînement de fluide extérieur à travers la frontière turbulente.}\label{fig:entrainement}
\end{figure}

\section{Le système résolu}
\[
\frac{\dd\theta}{\dd s}=\frac{c_f}{2}-\big(H+2\big)\frac{\theta}{U_e}\frac{\dd U_e}{\dd s},\qquad
\frac{\dd E}{\dd s}=U_e\,F(H_1),\qquad E=U_e\,\theta\,H_1,
\]
avec les fermetures empiriques suivantes (Head, puis Cebeci et Bradshaw) :
\[
F(H_1)=0{,}0306\,(H_1-3)^{-0{,}6169},
\]
\[
H_1(H)=\begin{cases}3{,}3+0{,}8234\,(H-1{,}1)^{-1{,}287} & H\le1{,}6,\\ 3{,}3+1{,}5501\,(H-0{,}6778)^{-3{,}064} & H>1{,}6,\end{cases}
\]
\[
c_f=0{,}246\cdot10^{-0{,}678\,H}\,\Re_\theta^{-0{,}268}\qquad\text{(Ludwieg-Tillmann).}
\]
Les inconnues intégrées sont $\theta$ et $E$. À chaque pas, on en déduit $H_1=E/(U_e\theta)$, puis $H$ en inversant $H_1(H)$, puis $c_f$.

\begin{figure}[htbp]
\centering
% figure générée par tools/make_figures.py : head_fermeture
\begin{tikzpicture}
\begin{groupplot}[group style={group size=2 by 1,horizontal sep=1.6cm},width=0.5\linewidth,height=5.2cm,
  title style={font=\small},tick label style={font=\scriptsize}]
\nextgroupplot[title={$H_1(H)$},xlabel={$H$},ylabel={$H_1=(\delta-\delta^*)/\theta$},ymin=3,ymax=12]
\addplot[cblue,very thick] coordinates {(1.15,42.208) (1.1639,31.694) (1.1777,25.35) (1.1916,21.151) (1.2055,18.19) (1.2193,16.001) (1.2332,14.326) (1.2471,13.006) (1.2609,11.943) (1.2748,11.071) (1.2887,10.344) (1.3025,9.7296) (1.3164,9.2043) (1.3303,8.7507) (1.3441,8.3556) (1.358,8.0086) (1.3718,7.7018) (1.3857,7.4288) (1.3996,7.1845) (1.4134,6.9648) (1.4273,6.7662) (1.4412,6.586) (1.455,6.4218) (1.4689,6.2716) (1.4828,6.1338) (1.4966,6.0069) (1.5105,5.8898) (1.5244,5.7814) (1.5382,5.6809) (1.5521,5.5873) (1.566,5.5001) (1.5798,5.4186) (1.5937,5.3423) (1.6076,5.2376) (1.6214,5.1517) (1.6353,5.0708) (1.6492,4.9944) (1.663,4.9224) (1.6769,4.8544) (1.6908,4.7901) (1.7046,4.7294) (1.7185,4.6718) (1.7324,4.6173) (1.7462,4.5656) (1.7601,4.5166) (1.7739,4.47) (1.7878,4.4258) (1.8017,4.3838) (1.8155,4.3438) (1.8294,4.3058) (1.8433,4.2696) (1.8571,4.2351) (1.871,4.2022) (1.8849,4.1708) (1.8987,4.1409) (1.9126,4.1123) (1.9265,4.085) (1.9403,4.0588) (1.9542,4.0339) (1.9681,4.01) (1.9819,3.9871) (1.9958,3.9652) (2.0097,3.9442) (2.0235,3.9241) (2.0374,3.9048) (2.0513,3.8863) (2.0651,3.8685) (2.079,3.8514) (2.0929,3.835) (2.1067,3.8193) (2.1206,3.8042) (2.1345,3.7896) (2.1483,3.7756) (2.1622,3.7621) (2.1761,3.7491) (2.1899,3.7366) (2.2038,3.7246) (2.2176,3.713) (2.2315,3.7018) (2.2454,3.691) (2.2592,3.6806) (2.2731,3.6705) (2.287,3.6609) (2.3008,3.6515) (2.3147,3.6424) (2.3286,3.6337) (2.3424,3.6253) (2.3563,3.6171) (2.3702,3.6092) (2.384,3.6016) (2.3979,3.5942) (2.4118,3.587) (2.4256,3.5801) (2.4395,3.5734) (2.4534,3.5669) (2.4672,3.5606) (2.4811,3.5545) (2.495,3.5486) (2.5088,3.5429) (2.5227,3.5374) (2.5366,3.532) (2.5504,3.5268) (2.5643,3.5217) (2.5782,3.5168) (2.592,3.512) (2.6059,3.5074) (2.6197,3.5029) (2.6336,3.4985) (2.6475,3.4942) (2.6613,3.4901) (2.6752,3.4861) (2.6891,3.4822) (2.7029,3.4784) (2.7168,3.4747) (2.7307,3.4711) (2.7445,3.4676) (2.7584,3.4642) (2.7723,3.4609) (2.7861,3.4577) (2.8,3.4546)};
\addplot[cmuted,dashed] coordinates {(2.4,3) (2.4,12)};
\node[cmuted,font=\scriptsize,anchor=south west] at (axis cs:2.42,8) {$H_{sep}=2{,}4$};
\nextgroupplot[title={entraînement $F(H_1)$},xlabel={$H_1$},ylabel={$F=\frac{1}{U_e}\frac{d(U_e\theta H_1)}{ds}$}]
\addplot[corange,very thick] coordinates {(3.35,0.058477) (3.4227,0.052051) (3.4954,0.047197) (3.5681,0.043374) (3.6408,0.040269) (3.7134,0.037686) (3.7861,0.035497) (3.8588,0.033612) (3.9315,0.031969) (4.0042,0.030521) (4.0769,0.029233) (4.1496,0.028079) (4.2223,0.027036) (4.295,0.02609) (4.3676,0.025225) (4.4403,0.024432) (4.513,0.023701) (4.5857,0.023025) (4.6584,0.022397) (4.7311,0.021812) (4.8038,0.021266) (4.8765,0.020754) (4.9492,0.020273) (5.0218,0.01982) (5.0945,0.019393) (5.1672,0.018989) (5.2399,0.018606) (5.3126,0.018243) (5.3853,0.017898) (5.458,0.01757) (5.5307,0.017257) (5.6034,0.016958) (5.6761,0.016672) (5.7487,0.016399) (5.8214,0.016137) (5.8941,0.015886) (5.9668,0.015645) (6.0395,0.015413) (6.1122,0.01519) (6.1849,0.014975) (6.2576,0.014768) (6.3303,0.014568) (6.4029,0.014375) (6.4756,0.014189) (6.5483,0.014009) (6.621,0.013835) (6.6937,0.013666) (6.7664,0.013503) (6.8391,0.013345) (6.9118,0.013191) (6.9845,0.013042) (7.0571,0.012898) (7.1298,0.012757) (7.2025,0.012621) (7.2752,0.012488) (7.3479,0.012359) (7.4206,0.012233) (7.4933,0.01211) (7.566,0.011991) (7.6387,0.011875) (7.7113,0.011761) (7.784,0.011651) (7.8567,0.011543) (7.9294,0.011438) (8.0021,0.011335) (8.0748,0.011234) (8.1475,0.011136) (8.2202,0.01104) (8.2929,0.010947) (8.3655,0.010855) (8.4382,0.010765) (8.5109,0.010677) (8.5836,0.010591) (8.6563,0.010507) (8.729,0.010425) (8.8017,0.010344) (8.8744,0.010265) (8.9471,0.010187) (9.0197,0.010111) (9.0924,0.010037) (9.1651,0.0099634) (9.2378,0.0098916) (9.3105,0.0098211) (9.3832,0.009752) (9.4559,0.0096841) (9.5286,0.0096175) (9.6013,0.009552) (9.6739,0.0094877) (9.7466,0.0094245) (9.8193,0.0093624) (9.892,0.0093013) (9.9647,0.0092413) (10.037,0.0091823) (10.11,0.0091243) (10.183,0.0090672) (10.255,0.0090111) (10.328,0.0089558) (10.401,0.0089015) (10.474,0.008848) (10.546,0.0087953) (10.619,0.0087434) (10.692,0.0086924) (10.764,0.0086421) (10.837,0.0085925) (10.91,0.0085437) (10.982,0.0084957) (11.055,0.0084483) (11.128,0.0084016) (11.2,0.0083556) (11.273,0.0083102) (11.346,0.0082655) (11.418,0.0082214) (11.491,0.0081779) (11.564,0.008135) (11.637,0.0080927) (11.709,0.008051) (11.782,0.0080098) (11.855,0.0079692) (11.927,0.0079291) (12,0.0078895)};
\end{groupplot}
\end{tikzpicture}

\caption{Les fermetures de Head. $H_1$ décroît quand $H$ croît : un profil creusé ($H$ grand) a peu de débit pour son déficit de quantité de mouvement. L'entraînement $F$ décroît avec $H_1$.}\label{fig:head-fermeture}
\end{figure}

\paragraph{Mécanisme du décollement turbulent.} En gradient adverse ($U_e'<0$) :
\begin{enumerate}
\item le terme $-(H+2)\theta U_e'/U_e$ devient positif et $\theta$ croît vite ;
\item $E$ croît plus lentement que $U_e\theta$, donc $H_1$ diminue et $H$ augmente : le profil se creuse ;
\item $c_f$ diminue avec $H$ (facteur $10^{-0{,}678H}$) ;
\item quand $H$ atteint 2 à 2,8, le frottement s'annule : c'est le décollement.
\end{enumerate}
Le seuil est flou dans la littérature. Le code prend $H_{sep}=2{,}4$ (constante \code{H\_SEP\_TURB}).

\section{Conditions initiales à la transition}
À la transition, on suppose $\theta$ continu (la quantité de mouvement ne saute pas) et on impose $H=1{,}4$, valeur typique d'une couche turbulente jeune. Le facteur de forme saute donc d'environ 2,6 (laminaire) à 1,4 (turbulent) : c'est le saut visible sur la figure~\ref{fig:bl-2412}. On démarre en $s_{tr}$ interpolé, avec $E_0=U_e(s_{tr})\,\theta(s_{tr})\,H_1(1{,}4)$, où $H_1(1{,}4)=7{,}18$.

\begin{codebox}[ : la fonction \code{rhs} et son intégration]
\begin{lstlisting}
def rhs(si, y):
    th, E = y                          # E = Ue * theta * H1
    ue = Ue_f(si)                      # interpolation lineaire de Ue(s)
    H1 = max(E / (ue * th), H1_sep)    # au-dela du decollement : H gele a 2.4
    Hh = float(H_of_H1(H1))
    cft = float(cf_ludwieg_tillmann(Hh, ue * th * Re))
    dth = 0.5 * cft - (Hh + 2) * th / ue * dUe_f(si)
    dE = ue * 0.0306 * (H1 - 3.0) ** -0.6169
    return [dth, dE]

sol = solve_ivp(rhs, (s_tr, st[-1]), y0, t_eval=st, method="RK45", rtol=1e-6, atol=1e-12)
\end{lstlisting}
Deux protections rendent l'intégration robuste. D'abord, $H_1$ est borné par le bas à $H_1(H_{sep})$ : au-delà du décollement, $H$ est gelé à 2,4 au lieu de diverger, et le drapeau \code{turbulent\_separation} est levé avec l'abscisse $x_{sep}$. Ensuite, \code{H\_of\_H1} borne $H_1$ au-dessus de 3,3, car l'inversion de $H_1(H)$ est singulière en $H_1\to3{,}3$.
\end{codebox}

\section{Le résultat sur un profil}
\begin{figure}[htbp]
\centering
% figure générée par tools/make_figures.py : cl_2412
\begin{tikzpicture}
\begin{groupplot}[group style={group size=3 by 1,horizontal sep=1.2cm},width=0.35\linewidth,height=5cm,
  xlabel={$x/c$},xmin=0,xmax=1,title style={font=\small},tick label style={font=\scriptsize},
  legend style={draw=none,font=\scriptsize}]
\nextgroupplot[title={$10^3\,\theta/c$},legend pos=north west]
\addplot[cblue,very thick] coordinates {(0.0038514,0.025409) (0.0060621,0.028835) (0.0087685,0.03276) (0.011969,0.037061) (0.01566,0.041645) (0.01984,0.046438) (0.024505,0.051386) (0.029651,0.056445) (0.035273,0.061585) (0.041368,0.066783) (0.047928,0.072022) (0.054949,0.077292) (0.062423,0.082584) (0.070344,0.087895) (0.078703,0.093223) (0.087494,0.098568) (0.096707,0.10393) (0.10633,0.10932) (0.11636,0.11473) (0.12679,0.12018) (0.13759,0.12565) (0.14877,0.13118) (0.16031,0.13675) (0.17219,0.14237) (0.18442,0.14805) (0.19696,0.1538) (0.20982,0.15963) (0.22297,0.16553) (0.2364,0.17153) (0.25011,0.17762) (0.26407,0.18382) (0.27826,0.19014) (0.29268,0.22309) (0.30731,0.25758) (0.32214,0.29173) (0.33714,0.32588) (0.3523,0.3603) (0.36761,0.39523) (0.38305,0.43102) (0.39857,0.46781) (0.41414,0.50511) (0.42974,0.54227) (0.44542,0.57929) (0.46115,0.61635) (0.47691,0.6536) (0.49269,0.69108) (0.50848,0.72884) (0.52425,0.7669) (0.54,0.8053) (0.5557,0.84403) (0.57135,0.88312) (0.58692,0.92258) (0.6024,0.96242) (0.61777,1.0027) (0.63302,1.0433) (0.64814,1.0844) (0.6631,1.126) (0.6779,1.1681) (0.69252,1.2107) (0.70695,1.2538) (0.72116,1.2976) (0.73515,1.3421) (0.74891,1.3873) (0.76242,1.4333) (0.77566,1.4802) (0.78863,1.5281) (0.80131,1.5771) (0.81369,1.6272) (0.82576,1.6787) (0.8375,1.7316) (0.84891,1.7861) (0.85997,1.8425) (0.87068,1.9009) (0.88102,1.9615) (0.89098,2.0247) (0.90055,2.0906) (0.90973,2.1599) (0.9185,2.2327) (0.92686,2.3098) (0.9348,2.3917) (0.94232,2.4792) (0.9494,2.5733) (0.95603,2.6751) (0.96222,2.786) (0.96795,2.908) (0.97323,3.0436) (0.97804,3.1965) (0.98238,3.3621)};\addlegendentry{extrados}
\addplot[corange,very thick] coordinates {(0.0040239,0.025927) (0.006247,0.025829) (0.0089501,0.028018) (0.01213,0.031211) (0.015782,0.034992) (0.019902,0.039175) (0.024485,0.04366) (0.029527,0.048382) (0.035021,0.053298) (0.040961,0.058375) (0.047342,0.063589) (0.054155,0.068922) (0.061395,0.074361) (0.069053,0.079893) (0.077121,0.085512) (0.085592,0.091209) (0.094457,0.09698) (0.10371,0.10282) (0.11333,0.10872) (0.12333,0.11469) (0.13367,0.12072) (0.14437,0.1268) (0.1554,0.13294) (0.16676,0.13912) (0.17844,0.14535) (0.19042,0.15162) (0.20269,0.15793) (0.21525,0.16427) (0.22808,0.17063) (0.24116,0.17701) (0.25449,0.1834) (0.26806,0.18979) (0.28185,0.19616) (0.29585,0.20251) (0.31005,0.20882) (0.32443,0.21507) (0.33899,0.22123) (0.3537,0.22728) (0.36855,0.23316) (0.38354,0.2388) (0.39866,0.244) (0.41396,0.24893) (0.42937,0.2542) (0.44486,0.25975) (0.4604,0.26543) (0.47599,0.27118) (0.4916,0.30364) (0.50723,0.33927) (0.52285,0.37397) (0.53845,0.40795) (0.55402,0.44136) (0.56954,0.47429) (0.58499,0.50677) (0.60037,0.53886) (0.61565,0.57057) (0.63082,0.60192) (0.64587,0.63291) (0.66078,0.66356) (0.67553,0.69387) (0.69012,0.72383) (0.70452,0.75346) (0.71872,0.78276) (0.73272,0.81175) (0.74649,0.84042) (0.76001,0.86881) (0.77329,0.89693) (0.7863,0.92478) (0.79903,0.9524) (0.81147,0.97984) (0.82361,1.0071) (0.83542,1.0343) (0.84691,1.0614) (0.85806,1.0885) (0.86886,1.1157) (0.8793,1.143) (0.88936,1.1706) (0.89904,1.1985) (0.90832,1.2269) (0.9172,1.2559) (0.92567,1.2857) (0.93372,1.3166) (0.94134,1.3488) (0.94852,1.3826) (0.95526,1.4184) (0.96155,1.4568) (0.96737,1.4984) (0.97274,1.5441) (0.97763,1.5953) (0.98205,1.6499)};\addlegendentry{intrados}
\nextgroupplot[title={facteur de forme $H$},ymin=1,ymax=3]
\addplot[cblue,very thick] coordinates {(0.0038514,2.4339) (0.0060621,2.4605) (0.0087685,2.4851) (0.011969,2.507) (0.01566,2.526) (0.01984,2.5422) (0.024505,2.5557) (0.029651,2.5671) (0.035273,2.5765) (0.041368,2.5845) (0.047928,2.5913) (0.054949,2.5972) (0.062423,2.6024) (0.070344,2.6073) (0.078703,2.6121) (0.087494,2.6167) (0.096707,2.6214) (0.10633,2.6263) (0.11636,2.6316) (0.12679,2.6373) (0.13759,2.6436) (0.14877,2.6505) (0.16031,2.6582) (0.17219,2.6667) (0.18442,2.6763) (0.19696,2.687) (0.20982,2.6991) (0.22297,2.7128) (0.2364,2.7284) (0.25011,2.7462) (0.26407,2.7668) (0.27826,2.7907) (0.29268,1.4161) (0.30731,1.4263) (0.32214,1.4321) (0.33714,1.4354) (0.3523,1.4375) (0.36761,1.4389) (0.38305,1.4403) (0.39857,1.4417) (0.41414,1.443) (0.42974,1.4438) (0.44542,1.4441) (0.46115,1.4442) (0.47691,1.4442) (0.49269,1.4443) (0.50848,1.4444) (0.52425,1.4446) (0.54,1.4448) (0.5557,1.4452) (0.57135,1.4457) (0.58692,1.4463) (0.6024,1.4471) (0.61777,1.4479) (0.63302,1.4489) (0.64814,1.4501) (0.6631,1.4514) (0.6779,1.4528) (0.69252,1.4544) (0.70695,1.4562) (0.72116,1.4582) (0.73515,1.4604) (0.74891,1.4629) (0.76242,1.4655) (0.77566,1.4684) (0.78863,1.4716) (0.80131,1.4752) (0.81369,1.4791) (0.82576,1.4833) (0.8375,1.488) (0.84891,1.4931) (0.85997,1.4988) (0.87068,1.5051) (0.88102,1.512) (0.89098,1.5197) (0.90055,1.5282) (0.90973,1.5378) (0.9185,1.5485) (0.92686,1.5607) (0.9348,1.5744) (0.94232,1.5903) (0.9494,1.6032) (0.95603,1.6194) (0.96222,1.6381) (0.96795,1.66) (0.97323,1.6862) (0.97804,1.7182) (0.98238,1.7548)};
\addplot[corange,very thick] coordinates {(0.0040239,2.2874) (0.006247,2.3541) (0.0089501,2.3887) (0.01213,2.412) (0.015782,2.431) (0.019902,2.4477) (0.024485,2.4628) (0.029527,2.4767) (0.035021,2.4894) (0.040961,2.5012) (0.047342,2.5121) (0.054155,2.5223) (0.061395,2.5319) (0.069053,2.541) (0.077121,2.5497) (0.085592,2.5581) (0.094457,2.5662) (0.10371,2.5741) (0.11333,2.5818) (0.12333,2.5894) (0.13367,2.5969) (0.14437,2.6042) (0.1554,2.6117) (0.16676,2.6189) (0.17844,2.6262) (0.19042,2.6336) (0.20269,2.6409) (0.21525,2.6481) (0.22808,2.6552) (0.24116,2.6621) (0.25449,2.6686) (0.26806,2.6746) (0.28185,2.6799) (0.29585,2.6844) (0.31005,2.6876) (0.32443,2.6893) (0.33899,2.6887) (0.3537,2.6851) (0.36855,2.6762) (0.38354,2.657) (0.39866,2.6364) (0.41396,2.6416) (0.42937,2.6693) (0.44486,2.6931) (0.4604,2.7101) (0.47599,2.7245) (0.4916,1.4106) (0.50723,1.4185) (0.52285,1.4234) (0.53845,1.4263) (0.55402,1.428) (0.56954,1.4289) (0.58499,1.4292) (0.60037,1.4293) (0.61565,1.4291) (0.63082,1.4287) (0.64587,1.4283) (0.66078,1.4279) (0.67553,1.4274) (0.69012,1.427) (0.70452,1.4266) (0.71872,1.4262) (0.73272,1.4259) (0.74649,1.4257) (0.76001,1.4255) (0.77329,1.4254) (0.7863,1.4255) (0.79903,1.4256) (0.81147,1.4258) (0.82361,1.4262) (0.83542,1.4268) (0.84691,1.4275) (0.85806,1.4283) (0.86886,1.4294) (0.8793,1.4307) (0.88936,1.4322) (0.89904,1.434) (0.90832,1.4361) (0.9172,1.4386) (0.92567,1.4415) (0.93372,1.4448) (0.94134,1.4487) (0.94852,1.4533) (0.95526,1.4587) (0.96155,1.4649) (0.96737,1.4724) (0.97274,1.4814) (0.97763,1.4924) (0.98205,1.5045)};
\nextgroupplot[title={$10^3\,c_f$},ymin=0,ymax=8]
\addplot[cblue,very thick] coordinates {(0.0038514,6.941) (0.0060621,5.5144) (0.0087685,4.4794) (0.011969,3.7164) (0.01566,3.1445) (0.01984,2.7081) (0.024505,2.369) (0.029651,2.1007) (0.035273,1.8848) (0.041368,1.708) (0.047928,1.5609) (0.054949,1.4367) (0.062423,1.3304) (0.070344,1.2382) (0.078703,1.1571) (0.087494,1.0852) (0.096707,1.0206) (0.10633,0.96212) (0.11636,0.90864) (0.12679,0.85934) (0.13759,0.81355) (0.14877,0.77073) (0.16031,0.73041) (0.17219,0.69221) (0.18442,0.6558) (0.19696,0.62089) (0.20982,0.58722) (0.22297,0.55455) (0.2364,0.52266) (0.25011,0.49131) (0.26407,0.46025) (0.27826,0.42919) (0.29268,4.3902) (0.30731,4.1622) (0.32214,3.994) (0.33714,3.8616) (0.3523,3.7517) (0.36761,3.6564) (0.38305,3.5699) (0.39857,3.4905) (0.41414,3.4186) (0.42974,3.3554) (0.44542,3.2999) (0.46115,3.2497) (0.47691,3.2034) (0.49269,3.1602) (0.50848,3.1193) (0.52425,3.0805) (0.54,3.0434) (0.5557,3.0077) (0.57135,2.9734) (0.58692,2.9401) (0.6024,2.9077) (0.61777,2.8761) (0.63302,2.8452) (0.64814,2.8149) (0.6631,2.7851) (0.6779,2.7556) (0.69252,2.7263) (0.70695,2.6972) (0.72116,2.6682) (0.73515,2.6391) (0.74891,2.6099) (0.76242,2.5805) (0.77566,2.5507) (0.78863,2.5204) (0.80131,2.4895) (0.81369,2.4579) (0.82576,2.4255) (0.8375,2.392) (0.84891,2.3574) (0.85997,2.3215) (0.87068,2.2841) (0.88102,2.2449) (0.89098,2.2037) (0.90055,2.1603) (0.90973,2.1142) (0.9185,2.0652) (0.92686,2.0126) (0.9348,1.9561) (0.94232,1.8947) (0.9494,1.8431) (0.95603,1.7833) (0.96222,1.7181) (0.96795,1.6462) (0.97323,1.5661) (0.97804,1.4755) (0.98238,1.3797)};
\addplot[corange,very thick] coordinates {(0.0040239,25.307) (0.006247,16.678) (0.0089501,11.981) (0.01213,9.112) (0.015782,7.2118) (0.019902,5.8823) (0.024485,4.9133) (0.029527,4.1837) (0.035021,3.6194) (0.040961,3.173) (0.047342,2.813) (0.054155,2.5176) (0.061395,2.2717) (0.069053,2.0642) (0.077121,1.8871) (0.085592,1.7343) (0.094457,1.6013) (0.10371,1.4845) (0.11333,1.3811) (0.12333,1.2891) (0.13367,1.2066) (0.14437,1.1324) (0.1554,1.0652) (0.16676,1.0042) (0.17844,0.94858) (0.19042,0.89781) (0.20269,0.85136) (0.21525,0.80883) (0.22808,0.76987) (0.24116,0.73423) (0.25449,0.7017) (0.26806,0.67213) (0.28185,0.64542) (0.29585,0.62153) (0.31005,0.60051) (0.32443,0.58251) (0.33899,0.56787) (0.3537,0.55731) (0.36855,0.55263) (0.38354,0.55915) (0.39866,0.56867) (0.41396,0.55229) (0.42937,0.51493) (0.44486,0.4839) (0.4604,0.46051) (0.47599,0.44056) (0.4916,4.3816) (0.50723,4.2031) (0.52285,4.0659) (0.53845,3.9562) (0.55402,3.8655) (0.56954,3.7884) (0.58499,3.7217) (0.60037,3.6628) (0.61565,3.6102) (0.63082,3.5626) (0.64587,3.5192) (0.66078,3.4792) (0.67553,3.4421) (0.69012,3.4075) (0.70452,3.3751) (0.71872,3.3445) (0.73272,3.3155) (0.74649,3.2878) (0.76001,3.2613) (0.77329,3.2358) (0.7863,3.2112) (0.79903,3.1872) (0.81147,3.1638) (0.82361,3.1407) (0.83542,3.1178) (0.84691,3.095) (0.85806,3.0722) (0.86886,3.0492) (0.8793,3.0258) (0.88936,3.0018) (0.89904,2.9771) (0.90832,2.9515) (0.9172,2.9247) (0.92567,2.8965) (0.93372,2.8664) (0.94134,2.834) (0.94852,2.799) (0.95526,2.7607) (0.96155,2.7185) (0.96737,2.6714) (0.97274,2.6181) (0.97763,2.5566) (0.98205,2.4924)};
\end{groupplot}
\end{tikzpicture}

\caption{Couche limite complète du NACA 2412 à $\alpha=2^\circ$, $\Re=3\cdot10^6$, calculée par \code{march\_surface}. On voit la transition (saut de $H$ et de $c_f$) en $x/c=0{,}28$ sur l'extrados et $0{,}48$ sur l'intrados, puis la montée de $H$ en fin d'extrados sous l'effet du gradient adverse.}\label{fig:bl-2412}
\end{figure}

La validation sur plaque plane turbulente (figure~\ref{fig:plaque}, à droite) donne $\theta$ à moins de 5~\% de la loi en $1/7$, $\theta=0{,}036\,x\,\Re_x^{-1/5}$, et $H\to1{,}34$.

\begin{piege}
\textbf{Le faux décollement de bord de fuite.} Près d'un bord de fuite fermé, l'écoulement potentiel décélère fortement (ch.~\ref{ch:trainee}), et ce gradient adverse est exagéré par rapport à l'écoulement réel. Sur la base de 512 profils à $\alpha=2^\circ$, 145 profils épais ou cambrés atteignent $H>2{,}4$ entre $x/c=0{,}957$ et $0{,}983$. Les déclarer invalides retirerait 28~\% de l'espace de conception à cause d'un artefact local. Le code les tolère (seuil \code{x\_sep\_crit}, ch.~\ref{ch:modele}).
\end{piege}

\begin{piege}
Head est une méthode « d'équilibre » : $F$ ne dépend que de la forme locale. Elle réagit donc instantanément, donc mal, aux changements brusques (choc, recollement de bulbe). Les méthodes à retard (\emph{lag-entrainment} de Green, ou le \emph{shear-lag} d'XFOIL) ajoutent une équation pour la contrainte de cisaillement et corrigent ce défaut.
\end{piege}

\exercices
\begin{exercice}
Calculer $H_1$ pour $H=1{,}3$, $1{,}4$ et $2{,}4$ avec les corrélations de Head. Que vaut $F$ pour $H=1{,}4$ ? Pour une couche limite d'épaisseur de quantité de mouvement $\theta=10^{-3}c$, quelle est alors l'épaisseur $\delta-\delta^*$ ?
\end{exercice}
\begin{exercice}
Montrer, à partir du système de Head et du sens de variation des fermetures, qu'en gradient adverse $H$ augmente. On pourra raisonner à $\theta$ fixé sur une petite distance et regarder comment évolue $H_1=E/(U_e\theta)$ quand $U_e$ diminue.
\end{exercice}
